\section{Conclusões}
\label{sec:conclusion}

Este trabalho aplicou os instrumentos clássicos da teoria amostral
--- fórmula de Cochran (1977), alocação ótima de Neyman e estimadores tipo
razão e tipo regressão sob desenho estratificado --- a um \textit{corpus}
público de toxicidade em português brasileiro (ToLD-Br), com o objetivo de
estabelecer um \textit{pipeline} reproduzível de auditoria estatística sobre
\textit{datasets} de detecção de dano.

\paragraph{Resposta às questões orientadoras.} \textbf{(Q1)} A fórmula de
Cochran, com correção de população finita, fornece $n=378$ como tamanho
amostral mínimo para estimar parâmetros populacionais com erro de amostragem
de $5\%$ a $95\%$ de confiança sobre o ToLD-Br. Este número é diretamente
verificável e dispensa o uso de amostras por conveniência. \textbf{(Q2)} A
alocação ótima de Neyman reduz o erro-padrão das estimativas em até $20\%$
frente à alocação proporcional para o mesmo orçamento $n$, com magnitude
dependente do estimador adotado. \textbf{(Q3)} Entre os cinco estimadores
avaliados, o tipo razão combinado apresentou o menor erro absoluto contra a
verdade populacional ($0{,}69\%$ de $\bar{Y}$); o tipo regressão separado
apresentou o menor erro-padrão. A baixa correlação observada entre a variável
auxiliar (\texttt{word\_count}) e a variável-alvo (\texttt{total\_severity})
limita o ganho da regressão sobre a média estratificada simples ---
limitação que sugere o uso de auxiliares mais informativos em trabalhos
futuros.

\paragraph{Contribuição à teoria.} Este trabalho fornece uma aplicação
explícita de Cochran §5.5--§7.11 e Bolfarine \& Bussab §5--§7 a um
\textit{corpus} de PNL em português brasileiro, com derivação completa das
variâncias amostrais sob FPC e validação empírica via $74$ testes unitários
contra propriedades textbook (recuperação exata sob censo total; redução de
variância da regressão face à média sob $y = 2x + \epsilon$).

\paragraph{Contribuição à técnica.} Disponibiliza-se um módulo Python
\textit{open source} (\texttt{toxicity\_analysis}) que encapsula o
\textit{pipeline} de amostragem, cálculo de \textit{features} léxicas e
estimação estratificada, com gerenciamento de dependências via \texttt{uv} e
amparo em $74$ testes \texttt{pytest}. O artefato é prontamente reutilizável
em outros \textit{datasets} de detecção de toxicidade.

\paragraph{Contribuição à sociedade.} A auditoria de modelos de linguagem é
uma componente crescente da regulação e da prática responsável de IA. Ao
fornecer instrumentos estatísticos calibrados para estimar prevalências e
intensidades de dano sobre \textit{corpora} de referência --- com erros de
amostragem explícitos --- o trabalho contribui para que decisões de
moderação automática e auditoria regulatória se apoiem em base inferencial
defensável.

\paragraph{Trabalhos futuros.}
(i) Extensão do \textit{pipeline} a outros \textit{corpora} em português
(HateBR, OLID-BR) e a \textit{datasets} bilíngues, comparando o ganho de
Neyman entre dispersões linguísticas distintas.
(ii) Investigação do esquema ordinal de votos como dimensão adicional de
estratificação --- por exemplo, por \textit{ratings} de concordância --- em
substituição à binarização por maioria.
(iii) Auditoria sobre saídas de LLMs (não apenas \textit{corpora} rotulados),
aplicando o mesmo desenho estratificado a amostras de gerações de modelos
treinados.
(iv) Avaliação sistemática de variáveis auxiliares mais ricas
(\textit{embeddings}, \texttt{caps\_ratio}, \texttt{punc\_density}) para
maximizar o ganho do estimador regressão.

\subsection*{Considerações Éticas}

Os \textit{datasets} de toxicidade utilizados como instrumentos de auditoria
não são neutros: refletem os critérios de rotulagem e os vieses culturais
dos próprios anotadores. O uso de tais \textit{corpora} para validar modelos
em escala carrega o risco de internalizar como objetivo um proxy ruidoso da
norma social vigente em um momento e contexto específicos. Recomenda-se que
auditorias estatísticas, ainda que metodologicamente rigorosas como a aqui
proposta, sejam acompanhadas de discussão qualitativa sobre os critérios
implícitos no \textit{corpus}-fonte e suas limitações de cobertura cultural.

\subsection*{Disponibilidade e Reprodutibilidade}

Código-fonte, \textit{notebooks}, suíte de testes e \textit{pipeline} de
\textit{download} do ToLD-Br estão disponíveis em
\texttt{https://github.com/[autor]/toxicity-analysis} (a ser publicado).
